\section{Compito 1}

\subsection{Descrizione del Problema}
\subsubsection{Descrizione informale del problema}
È dato un ambiente discreto rappresentato da una griglia quadrata $N \times N$.
Ogni cella è identificata dalla coppia di coordinate $(r, c)$ con $0 \le r, c \le N-1$ ed è di uno di due tipi:
\begin{itemize}
    \item \textbf{Libera}: il robot può occuparla;
    \item \textbf{Ostacolo}: il robot non può occuparla.
\end{itemize}

\noindent
Nella griglia sono individuate due celle libere particolari: la cella di partenza $S$ (\textit{Start}) e la cella obiettivo $G$ (\textit{Goal}).

\noindent\\
Da ogni cella libera il robot può tentare uno spostamento in una delle quattro direzioni cardinali:
\[
\mathcal{A} = \{\text{NORD},\ \text{SUD},\ \text{EST},\ \text{OVEST}\}.
\]
Lo spostamento avviene se la cella di arrivo è interna alla griglia ed è libera, altrimenti il robot resta fermo nella cella corrente.

\begin{table}[H]
\centering
\begin{tabular}{lc}
\hline
\textbf{Evento} & \textbf{Premio $R$} \\
\hline
Passo su una cella libera & $-1$ \\
Urto contro un ostacolo o contro il bordo & $-10$ \\
Raggiungimento della cella Goal & $+100$ \\
\hline
\end{tabular}
\caption{Struttura dei premi del modello.}
\end{table}

\noindent
\textbf{Obiettivo:} Determinare, per ogni cella libera della griglia, il valore della soluzione ottima $V^*$ e ricostruire un cammino ottimo da $S$ a $G$ (grazie alla matrice $\pi^*$).

\clearpage
\subsubsection{Formalizzazione}
Il problema viene formalizzato come un Processo Decisionale di Markov (MDP) deterministico a tempo discreto definito dalla tupla $(\mathcal{S}, \mathcal{A}, \delta, R, \gamma)$, unitamente alle funzioni di valore ($V^*, Q^*$) e alla politica ottima ($\pi^*$):

\begin{itemize}
    \item \textbf{Spazio degli stati} $\mathcal{S}$:
    L'insieme di tutte le celle libere della griglia $N \times N$:
    \[
    \mathcal{S} = \{(r,c) : 0 \le r,c \le N-1,\ \text{Grid}[r,c] \ne \text{OSTACOLO}\}
    \]
    La cella obiettivo $G \in \mathcal{S}$ è definita come \textbf{stato terminale assorbente}. Si noti che in generale $|\mathcal{S}| \le N^2$, con uguaglianza solo in assenza di ostacoli.

    \item \textbf{Spazio delle azioni} $\mathcal{A}$:
    L'insieme delle 4 direzioni cardinali ammissibili per il movimento del robot:
    \[
    \mathcal{A} = \{\text{NORD},\ \text{SUD},\ \text{EST},\ \text{OVEST}\}
    \]

    \item \textbf{Funzione di transizione} $\delta : \mathcal{S} \times \mathcal{A} \to \mathcal{S}$:
    Determina lo stato di destinazione in modo deterministico:
    \[
    \delta(s,a) = \begin{cases}
    s' & \text{se la cella adiacente in direzione } a \text{ è libera e dentro la griglia} \\
    s & \text{in caso di urto contro un ostacolo o il bordo}
    \end{cases}
    \]

    \item \textbf{Funzione di premio} $R : \mathcal{S} \times \mathcal{A} \to \mathbb{R}$:
    Rappresenta la ricompensa immediata ottenuta eseguendo l'azione $a$ dallo stato $s$:
    \[
    R(s,a) = \begin{cases}
    c_{\text{urto}} = -10 & \text{se } \delta(s,a) = s \text{ (urto)} \\
    c_{\text{raggiungimento}} = +100 & \text{se } \delta(s,a) = G \text{ (obiettivo raggiunto)} \\
    c_{\text{passo}} = -1 & \text{altrimenti (spostamento valido)}
    \end{cases}
    \]
    \textit{Nota}: Il costo $c_{\text{passo}} = -1$ incentiva il robot a trovare il cammino di lunghezza minima verso $G$.

    \item \textbf{Fattore di sconto} $\gamma = 1$: i premi futuri vengono valutati esattamente quanto quelli immediati, senza alcuna penalizzazione temporale. In combinazione con la struttura dei premi ciò garantisce che il cammino ottimo sia quello di lunghezza minima: percorrendo $L$ mosse da $s$ fino a $G$ il ritorno complessivo vale
    \[
    -(L-1) + 100 = 101 - L,
    \]
    quindi massimizzare il ritorno equivale esattamente a minimizzare $L$. Questa scelta ha però una conseguenza importante sulla terminazione dell'algoritmo, discussa nella sezione seguente.
\end{itemize}

\subsubsection{Ipotesi di raggiungibilità e terminazione}
Con $\gamma < 1$ l'operatore di Bellman è una contrazione di fattore $\gamma$ e la convergenza della Value Iteration è garantita per qualunque istanza. Con $\gamma = 1$ questa garanzia viene meno ed è necessaria un'ipotesi esplicita sul dominio:

\begin{quote}
\textbf{Ipotesi (R).} Da ogni cella libera $s \in \mathcal{S}$ esiste una sequenza di azioni che conduce alla cella obiettivo $G$.
\end{quote}

\noindent
L'ipotesi è necessaria: se una cella libera $s$ non potesse raggiungere $G$, ogni azione disponibile in $s$ produrrebbe un premio negativo ($-1$ oppure $-10$) senza mai incassare il premio terminale $+100$, e si avrebbe $V^*(s) = -\infty$. In quel caso i valori diminuirebbero indefinitamente a ogni iterazione, $\Delta$ non scenderebbe mai sotto $\epsilon$ e il ciclo \texttt{while} non terminerebbe.

\noindent\\
Sotto l'ipotesi (R) la convergenza è invece garantita e, come si dimostrerà nell'analisi asintotica, avviene in un numero finito e limitato di iterazioni. Poiché l'ipotesi non è verificabile a priori su un'istanza arbitraria, entrambi gli pseudocodici presentati includono un limite di sicurezza sul numero di iterazioni che ne garantisce comunque la terminazione, segnalando l'anomalia.

\clearpage
\noindent
Per risolvere l'MDP e determinare il cammino ottimo si definiscono tre entità fondamentali legate tra loro dalle Equazioni di Ottimalità di Bellman:

\begin{enumerate}
    \item \textbf{Funzione Valore Azione-Stato Ottima} $Q^* : \mathcal{S} \times \mathcal{A} \to \mathbb{R}$:
    Rappresenta il ritorno atteso eseguendo l'azione $a$ nello stato $s$ e proseguendo in modo ottimo da quel punto in poi:
    \[
    Q^*(s,a) = R(s,a) + \gamma \, V^*(\delta(s,a))
    \]

    \item \textbf{Funzione Valore Stato Ottima} $V^* : \mathcal{S} \to \mathbb{R}$:
    Rappresenta il massimo ritorno atteso scontato ottenibile a partire dallo stato $s$.
    \begin{itemize}
        \item Poiché $G$ è uno stato terminale assorbente da cui non si possono compiere ulteriori azioni, il suo valore è fisso:
        \[
        V^*(G) = 0
        \]
        \item Per ogni altro stato libero $s \in \mathcal{S} \setminus \{G\}$, soddisfa l'\textbf{Equazione di Ottimalità di Bellman}:
        \[
        V^*(s) = \max_{a \in \mathcal{A}} Q^*(s,a) = \max_{a \in \mathcal{A}} \left\{ R(s,a) + \gamma \, V^*(\delta(s,a)) \right\}
        \]
    \end{itemize}

    \item \textbf{Politica Ottima} $\pi^* : \mathcal{S} \setminus \{G\} \to \mathcal{A}$:
    È la funzione di decisione deterministica che associa a ogni stato l'azione ottima da compiere. Viene estratta a partire da $Q^*(s,a)$ mediante l'operatore $\arg\max$:
    \[
    \pi^*(s) = \arg\max_{a \in \mathcal{A}} Q^*(s,a) = \arg\max_{a \in \mathcal{A}} \left\{ R(s,a) + \gamma \, V^*(\delta(s,a)) \right\}
    \]
    Quando più azioni realizzano lo stesso valore massimo, l'$\arg\max$ non è univoco: la politica ottima non è quindi necessariamente unica, mentre $V^*$ lo è sempre.
\end{enumerate}

\clearpage
\subsection{Dimostrazione dell'applicabilità della Programmazione Dinamica}
\subsubsection{Primo Segno Distintivo: Sottostruttura Ottima}

Nel nostro problema:
\begin{itemize}
    \item Sia $P^* = (s_0, s_1, s_2, \dots, s_k = G)$ un cammino ottimo che conduce dalla cella di partenza $S = s_0$ alla cella obiettivo $G$.
    \item Se si considera una qualsiasi cella intermedia $s_i \in P^*$, il sottocammino $P^*_{i \to k} = (s_i, s_{i+1}, \dots, G)$ rappresenta necessariamente la soluzione ottima al sottoproblema di raggiungere $G$ a partire da $s_i$.
    \item Se infatti esistesse un sottocammino da $s_i$ a $G$ preferibile a $P^*_{i \to k}$, sostituendolo all'interno di $P^*$ si otterrebbe un cammino globale da $S$ a $G$ migliore di $P^*$, il che contraddirebbe l'ottimalità di $P^*$.
\end{itemize}

\noindent
Di conseguenza, la soluzione del problema principale per la cella $s$ si compone direttamente a partire dalla soluzione ottima del sottoproblema relativo allo stato adiacente $\delta(s, a^*)$ scelto dalla politica ottima.

\subsubsection{Secondo Segno Distintivo: Risoluzione dello stesso sottoproblema più volte}
Definizione ricorsiva della funzione valore $V(s)$ per una generica cella $s$:
\[
V(s) = \begin{cases}
0 & \text{se } s = G \\
\max_{a \in \mathcal{A}} \left\{ R(s,a) + \gamma \, V(\delta(s,a)) \right\} & \text{se } s \in \mathcal{S} \setminus \{G\}
\end{cases}
\]

\noindent
La definizione ricorsiva è difficile da scrivere come una ricorrenza in quanto né la dimensione della griglia diminuisce a ogni iterazione né diminuisce la distanza da $G$.
Si potrebbe esprimere la ricorsione in termini di orizzonte temporale, però abbiamo preferito lasciarla così in quanto è molto più intuitiva.\\
Importante notare che usando questo approccio ricorsivo senza l'uso della programmazione dinamica, molti sottoproblemi vengono risolti più volte, in quanto la funzione valore $V(s)$ viene calcolata per ogni cella libera della griglia, e ciascuna cella può essere raggiunta da più celle adiacenti.

\clearpage
\subsection{Pseudocodice dell'Algoritmo}
In questa sezione vengono presentati gli pseudocodici sviluppati secondo il paradigma della Programmazione Dinamica per la risoluzione del problema descritto. Il procedimento di calcolo si articola in due algoritmi distinti:

\begin{enumerate}
    \item \textbf{Calcolo del valore della soluzione ottima:} l'algoritmo \texttt{GRID-VALUE-ITERATION} adotta uno schema iterativo \textit{bottom-up} per calcolare la funzione valore ottima $V^*$ su tutta la griglia. Durante il calcolo, aggiorna contemporaneamente una tabella ausiliaria $\pi$ memorizzando le decisioni locali ottime.
    \item \textbf{Costruzione della soluzione ottima:} l'algoritmo \texttt{CONSTRUCT-OPTIMAL-PATH} sfrutta le informazioni salvate nella tabella $\pi$ dal primo algoritmo per tracciare e restituire il cammino ottimo sotto forma di sequenza ordinata di celle dalla partenza $S$ all'obiettivo $G$.
\end{enumerate}

\subsubsection{Pseudocodice per il valore di una soluzione ottima}
L'algoritmo \texttt{GRID-VALUE-ITERATION} riceve in ingresso la griglia, la cella obiettivo $G$, il fattore di sconto $\gamma$ e la tolleranza $\epsilon > 0$. Il ciclo principale \texttt{while} esegue aggiornamenti successivi della tabella dei valori $V$ fino a convergenza, ovvero fino a quando la massima variazione dei valori tra due iterazioni consecutive ($\Delta$) risulta inferiore a $\epsilon$. Il contatore $k$, confrontato con il limite $max\_iter = N^2 + 1$, garantisce la terminazione anche quando l'ipotesi (R) non è soddisfatta.

\paragraph{Ingresso:}
\begin{itemize}
    \item $Grid[0..N-1, 0..N-1]$: matrice $N \times N$ rappresentante l'ambiente ($Grid[r,c] = \text{LIBERA}$ oppure $\text{OSTACOLO}$);
    \item $N$: dimensione del lato della griglia quadrata;
    \item $G = (r_G, c_G)$: coordinate della cella obiettivo (\textit{Goal});
    \item $\gamma$: fattore di sconto;
    \item $\epsilon$: soglia di tolleranza per il criterio di convergenza ($\epsilon > 0$).
\end{itemize}

\paragraph{Uscite:}
\begin{itemize}
    \item $V[0..N-1, 0..N-1]$: tabella ausiliaria contenente i valori ottimi $V^*(s)$ per ogni stato $s = (r,c)$;
    \item $\pi[0..N-1, 0..N-1]$: tabella ausiliaria per la memorizzazione dell'azione ottima $\pi^*(s) \in \mathcal{A}$ da intraprendere da ciascuna cella $s = (r,c)$.
\end{itemize}

\noindent
Per leggibilità il calcolo del massimo su $\mathcal{A}$ per una singola cella è isolato nella procedura ausiliaria \texttt{BELLMAN-UPDATE} (Algoritmo~\ref{alg:bellman}), che applica l'Equazione di Ottimalità di Bellman alla cella $s$ leggendo i valori dell'iterazione precedente. Nel seguito, per $s = (r,c)$, si scrive $V[s]$ in luogo di $V[r,c]$ (analogamente per $V_{old}$, $\pi$ e $Grid$).

\begin{algorithm}[H]
\caption{Aggiornamento di Bellman per una singola cella}
\label{alg:bellman}
\begin{algorithmic}[1]
\Statex \textbf{BELLMAN-UPDATE}$(s, V_{old}, \gamma)$
\State $max\_q \gets -\infty$;\ \ $best\_a \gets \text{NIL}$
\For{$a \in \{\text{NORD}, \text{SUD}, \text{EST}, \text{OVEST}\}$} \label{bu:for}
    \State \Do $s' \gets \delta(s, a)$
    \State $q \gets R(s, a) + \gamma \cdot V_{old}[s']$
    \If{$q > max\_q$}
        \State \Then $max\_q \gets q$;\ \ $best\_a \gets a$ \label{bu:endfor}
    \EndIf
\EndFor
\State \Return $(max\_q, best\_a)$
\end{algorithmic}
\end{algorithm}

\begin{algorithm}[H]
\caption{Calcolo del valore di una soluzione ottima}
\label{alg:gvi}
\begin{algorithmic}[1]
\Statex \textbf{GRID-VALUE-ITERATION}$(Grid, N, G, \gamma, \epsilon)$
\For{$r \gets 0$ \textbf{to} $N-1$} \label{gvi:init-start}
    \DoFor{$c \gets 0$ \textbf{to} $N-1$}
        \State \Do $V[r, c] \gets 0$;\ \ $\pi[r, c] \gets \text{NIL}$
    \EndDoFor
\EndFor
\State $\Delta \gets \infty$;\ \ $k \gets 0$;\ \ $max\_iter \gets N^2 + 1$ \label{gvi:init-end}
\While{$\Delta \ge \epsilon$ \textbf{and} $k < max\_iter$} \label{gvi:while}
    \State \Do $\Delta \gets 0$;\ \ $k \gets k + 1$
    \State $V_{old} \gets V$ \label{gvi:copy}
    \For{$r \gets 0$ \textbf{to} $N-1$} \label{gvi:scan}
        \DoFor{$c \gets 0$ \textbf{to} $N-1$} \label{gvi:scan-c}
            \State \Do $s \gets (r, c)$
            \If{$Grid[s] \ne \text{OSTACOLO}$ \textbf{and} $s \ne G$}
                \State \Then $(V[s], \pi[s]) \gets \text{BELLMAN-UPDATE}(s, V_{old}, \gamma)$
                \State $\Delta \gets \max(\Delta, |V_{old}[s] - V[s]|)$ \label{gvi:while-end}
            \EndIf
        \EndDoFor
    \EndFor
\EndWhile
\If{$\Delta \ge \epsilon$}
    \State \Then \textbf{error} ``Nessuna convergenza: esistono celle libere da cui $G$ è irraggiungibile''
\EndIf
\State \Return $(V, \pi)$
\end{algorithmic}
\end{algorithm}

\clearpage
\subsubsection{Pseudocodice per la costruzione di una soluzione ottima}
Calcolata la tabella delle decisioni ottime $\pi$ mediante \texttt{GRID-VALUE-ITERATION}, la ricostruzione della sequenza di celle da $S$ a $G$ non richiede ulteriori valutazioni della funzione di valore, ma un semplice attraversamento deterministico guidato dalla tabella $\pi$.

\noindent
L'algoritmo \texttt{CONSTRUCT-OPTIMAL-PATH} compone la sequenza ordinata di coordinate $P = \langle s_0, s_1, \dots, s_k \rangle$ con $s_0 = S$ e $s_k = G$. Per garantire la terminazione dell'algoritmo anche in presenza di griglie malformate o prive di soluzioni, viene introdotto un controllo sul numero massimo di passi eseguiti ($N^2$).

\begin{algorithm}[H]
\caption{Costruzione di una soluzione ottima}
\label{alg:path}
\begin{algorithmic}[1]
\Statex \textbf{CONSTRUCT-OPTIMAL-PATH}$(S, G, \pi, Grid, N)$
\State $s \gets S$;\ \ $P \gets \text{Queue()}$;\ \ $\text{Enqueue}(P, s)$;\ \ $steps \gets 0$ \label{cop:init}
\While{$s \ne G$ \textbf{and} $steps < N^2$} \label{cop:while}
    \State \Do $a^* \gets \pi[s]$
    \If{$a^* = \text{NIL}$ \textbf{or} $\delta(s, a^*) = s$}
        \State \Then \textbf{error} ``Mossa ottima non definita o diretta contro un ostacolo/bordo''
    \EndIf
    \State $s \gets \delta(s, a^*)$
    \State $\text{Enqueue}(P, s)$;\ \ $steps \gets steps + 1$ \label{cop:while-end}
\EndWhile
\If{$s \ne G$}
    \State \Then \textbf{error} ``Goal non raggiungibile entro il numero massimo di passi consentiti''
\EndIf
\State \Return $P$
\end{algorithmic}
\end{algorithm}


\subsection{Analisi Asintotica}
\subsubsection{Complessità Temporale}

\paragraph{Algoritmo \texttt{GRID-VALUE-ITERATION}:}
\begin{itemize}
    \item \textbf{Inizializzazione (righe \ref{gvi:init-start}--\ref{gvi:init-end}):} I due cicli nidificati sulle righe $r$ e colonne $c$ scorrono tutte le $N \times N$ celle della matrice per azzerare $V$ e impostare $\pi$ a $\text{NIL}$. Ciascuna assegnazione richiede tempo $\Theta(1)$. La complessità dell'inizializzazione è dunque $\Theta(N^2)$.
    \item \textbf{Ciclo principale di convergenza (righe \ref{gvi:while}--\ref{gvi:while-end}):} Sia $K$ il numero totale di iterazioni del ciclo \texttt{while} eseguite prima che la variazione massima rispetti il criterio di arresto $\Delta < \epsilon$. In ciascuna delle $K$ iterazioni vengono eseguite le seguenti operazioni:
    \begin{itemize}
        \item La copia $V_{old} \gets V$ (riga \ref{gvi:copy}) duplica l'intera matrice $N \times N$ e costa quindi $\Theta(N^2)$;
        \item Scansione completa della griglia mediante due cicli \texttt{for} nidificati (righe \ref{gvi:scan}--\ref{gvi:scan-c}) sulle $N \times N$ celle, per un totale di $N^2$ stati valutati;
        \item Per ogni cella ammissibile (non ostacolo e diversa da $G$), la chiamata a \texttt{BELLMAN-UPDATE} esegue il ciclo dell'Algoritmo~\ref{alg:bellman} (righe \ref{bu:for}--\ref{bu:endfor}), che valuta le $|\mathcal{A}| = 4$ direzioni cardinali. Il calcolo dello stato di destinazione $s' = \delta(s,a)$, la lettura di $R(s,a)$ e l'accesso in tabella a $V_{old}[s']$ avvengono in tempo costante $\Theta(1)$;
        \item L'aggiornamento dei valori $max\_q$, $best\_a$, $V[s]$, $\pi[s]$ e il calcolo del valore assoluto per $\Delta$ richiedono tempo $\Theta(1)$.
    \end{itemize}
    Il costo computazionale di una singola iterazione del ciclo \texttt{while} è quindi $\Theta(N^2) + \Theta(N^2 \cdot |\mathcal{A}|) = \Theta(N^2)$.
\end{itemize}

\noindent
Moltiplicando il costo di una singola iterazione per il numero complessivo di iterazioni $K$, la complessità temporale di \texttt{GRID-VALUE-ITERATION} risulta:
\[
T_{\text{Value-Iteration}}(N, K) = \Theta(N^2) + \Theta(K \cdot N^2) = \Theta(K \cdot N^2)
\]

\paragraph{Limitazione del numero di iterazioni $K$.}
Il parametro $K$ non è indipendente dall'istanza. Sia $d(s,G)$ la distanza minima da $s$ a $G$, misurata in numero di mosse sul grafo delle celle libere. Per $\gamma = 1$ e con inizializzazione $V_0 \equiv 0$ si dimostra per induzione che, dopo $k \ge 1$ iterazioni, per ogni cella $s \ne G$ vale
\[
V_k(s) = \begin{cases}
101 - d(s,G) = V^*(s) & \text{se } d(s,G) \le k \\
-k & \text{se } d(s,G) > k
\end{cases}
\]
Infatti $V_k(s)$ è il miglior ritorno ottenibile in al più $k$ mosse: se $G$ è raggiungibile entro $k$ mosse conviene raggiungerlo per la via più breve, altrimenti la scelta migliore è compiere $k$ passi su celle libere (premio $-1$ ciascuno), evitando gli urti ($-10$).

\noindent\\
Ne segue che l'iterazione $k$ modifica esattamente le celle con $d(s,G) \ge k$: quelle a distanza $k$ passano da $-(k-1)$ a $101-k$, quelle più lontane da $-(k-1)$ a $-k$. Detta $\text{ecc}(G) = \max_{s \in \mathcal{S}} d(s, G)$ l'\textbf{eccentricità} del Goal nel grafo delle celle libere, le iterazioni $1, \dots, \text{ecc}(G)$ producono quindi $\Delta \ge 1$, mentre l'iterazione $\text{ecc}(G)+1$ non modifica alcun valore e produce $\Delta = 0$, arrestando il ciclo. Vale quindi l'uguaglianza
\[
K_{\text{Standard}} = \text{ecc}(G) + 1
\]
e, poiché un cammino minimo non ripete stati, $\text{ecc}(G) \le |\mathcal{S}| - 1 \le N^2 - 1$, da cui
\[
K_{\text{Standard}} \le |\mathcal{S}| \le N^2 .
\]
Sostituendo nella formula precedente si ottiene la limitazione superiore per il caso peggiore:
\[
T_{\text{Value-Iteration}}(N) = O(N^2 \cdot N^2) = O(N^4)
\]
La limitazione vale anche quando l'ipotesi (R) non è soddisfatta, poiché in quel caso il ciclo viene interrotto dal limite di sicurezza dopo $max\_iter = N^2 + 1$ iterazioni. Sotto l'ipotesi (R), invece, $K_{\text{Standard}} \le N^2 < max\_iter$ e il limite di sicurezza non interviene mai.

\noindent\\
Il caso peggiore è raggiunto dalle istanze a corridoio unico (serpentine e labirinti a pettine), in cui l'unico cammino ammissibile attraversa una frazione costante delle celle libere e quindi $\text{ecc}(G) = \Theta(N^2)$. Nelle istanze poco ostruite si ha invece $\text{ecc}(G) = \Theta(N)$ e il costo scende a $\Theta(N^3)$. La verifica sperimentale di questa uguaglianza è riportata nel Compito 3.

\noindent\\
Si osservi infine che, essendo i premi interi e $\gamma = 1$, i valori calcolati sono interi e la convergenza avviene in modo \textit{esatto} ($\Delta$ diventa esattamente $0$): la tolleranza $\epsilon$ non influenza quindi il numero di iterazioni, purché $\epsilon \le 1$.

\clearpage
\paragraph{Algoritmo \texttt{CONSTRUCT-OPTIMAL-PATH}:}
\begin{itemize}
    \item \textbf{Inizializzazione (riga \ref{cop:init}):} L'assegnazione dello stato iniziale, la creazione della coda $P$ contenente $S$ e l'azzeramento del contatore richiedono tempo $\Theta(1)$.
    \item \textbf{Ciclo di ricostruzione del cammino (righe \ref{cop:while}--\ref{cop:while-end}):} Sia $L$ la lunghezza del percorso estratto (in numero di passi), con $0 \le L \le N^2$. Il ciclo \texttt{while} viene eseguito una volta per ciascun passo del cammino ottimo, cioè $L$ volte, e ogni iterazione aggiunge una cella a $P$:
    \begin{itemize}
        \item Ad ogni passo, l'accesso alla tabella della politica $\pi[s]$ avviene in tempo $\Theta(1)$;
        \item La funzione di transizione $\delta(s, a^*)$ e le operazioni di verifica sugli errori e sugli stalli richiedono tempo $\Theta(1)$;
        \item L'operazione $\text{Enqueue}(P, s)$ inserisce l'elemento in coda in tempo $\Theta(1)$.
    \end{itemize}
    Il ciclo termina quando $s = G$ (oppure se viene raggiunto il limite di sicurezza di $N^2$ passi).
\end{itemize}

\noindent
La complessità temporale di \texttt{CONSTRUCT-OPTIMAL-PATH} è quindi proporzionale alla lunghezza del cammino ottimo $L$:
\[
T_{\text{Construct-Path}}(L) = \Theta(L + 1)
\]
Nel caso peggiore in cui il cammino ottimo attraversi la quasi totalità delle celle libere della griglia allora $L = \Theta(N^2)$, mentre nel caso migliore $S = G$ e si ha $L = 0$. Di conseguenza possiamo affermare che la complessità temporale sia compresa, in generale, fra $\Omega(1)$ e $O(N^2)$.
\paragraph{Complessità Temporale Complessiva dell'intero procedimento:}
\[
T_{\text{Totale}}(N, K, L) = \Theta(K \cdot N^2) + \Theta(L + 1) = \Theta(K \cdot N^2) = O(N^4)
\]

\clearpage
\subsubsection{Complessità Spaziale}

\paragraph{1. Memoria allocata da \texttt{GRID-VALUE-ITERATION}:}
\begin{itemize}
    \item \textbf{Matrice della griglia $Grid[0..N-1, 0..N-1]$:} occupa una quantità di memoria pari a $\Theta(N^2)$ celle;
    \item \textbf{Tabella dei valori $V[0..N-1, 0..N-1]$:} matrice di numeri reali di dimensione $N \times N$, pari a $\Theta(N^2)$ spazio;
    \item \textbf{Tabella della politica $\pi[0..N-1, 0..N-1]$:} matrice di elementi dell'alfabeto $\mathcal{A} \cup \{\text{NIL}\}$ di dimensione $N \times N$, pari a $\Theta(N^2)$ spazio;
    \item \textbf{Copia di lavoro $V_{old}[0..N-1, 0..N-1]$ (riga \ref{gvi:copy}):} lo schema di aggiornamento \textbf{sincrono} richiede di conservare, per tutta la durata di ogni iterazione, l'intera tabella dei valori dell'iterazione precedente. Si tratta di una seconda matrice $N \times N$ di numeri reali, pari a ulteriori $\Theta(N^2)$ spazio. Questo è l'unico termine \textit{ausiliario} in senso stretto, cioè non richiesto dalla specifica del problema né dal formato dell'uscita: è precisamente il termine che la variante in-place del Compito 2 elimina;
    \item \textbf{Variabili scalari ausiliarie:} le variabili $s, \Delta, k, max\_iter$ e, in \texttt{BELLMAN-UPDATE}, $a, s', q, max\_q, best\_a$ occupano uno spazio costante $\Theta(1)$.
\end{itemize}

La complessità spaziale occupata dalle strutture tabulari è pertanto $\Theta(N^2)$, di cui $\Theta(N^2)$ di memoria \textit{ausiliaria} imputabile alla sola copia $V_{old}$.

\paragraph{2. Memoria allocata da \texttt{CONSTRUCT-OPTIMAL-PATH}:}
\begin{itemize}
    \item \textbf{Struttura dati Coda $P$:} memorizza la sequenza delle coordinate delle celle del cammino ottimo da $S$ a $G$. Nel caso di un cammino lungo $L$ passi, la coda contiene $L + 1$ coppie di coordinate $(r, c)$, occupando uno spazio pari a $\Theta(L + 1)$;
    \item \textbf{Variabili di controllo:} le variabili scalari $s, a^*, steps$ occupano memoria costante $\Theta(1)$.
\end{itemize}

\noindent
Poiché il percorso $P$ non può superare il numero totale di celle della griglia ($L \le N^2$), lo spazio massimo richiesto dalla coda è $O(N^2)$.

\paragraph{Complessità Spaziale Complessiva dell'intero procedimento:}
Considerando tutte le tabelle ausiliarie ($V$, $V_{old}$, $\pi$) e la struttura dati della soluzione ($P$), la memoria complessiva richiesta è dominata dalle matrici $N \times N$:
\[
S_{\text{Totale}}(N) = \Theta(N^2)
\]
